\documentclass[12pt]{article}
\usepackage[utf8]{inputenc}
\usepackage[T1]{fontenc}
\usepackage{hyperref}
\title{Warumungu Language Resources for K-12 Curriculum Development}
\author{Compiled by Claude Sonnet 5.5 (claude-sonnet-5-5) for Andr\'as Kornai}
\date{}
\begin{document}
\maketitle

\section{Introduction}
Warumungu (ISO 639-3 wrm; Glottocode waru1265; Pama-Nyungan) is spoken in and around Tennant Creek, Northern Territory. It has more published material than its neighbour Warlmanpa, including dictionaries and a learner's guide.

\section{Primary Educational Resources}
\begin{itemize}
\item \emph{Warumungu Dictionary} (1989, ed.\ Trisha Frank) and \emph{Warumungu Picture Dictionary} (1989, Edward Plummer) are digitised on Territory Stories under CC BY-NC-ND 4.0. Items from the Living Archive of Aboriginal Languages are at \url{https://territorystories.nt.gov.au/10070/882584} and \url{https://territorystories.nt.gov.au/10070/882586/0}. The results do not say which record is which.
\item \emph{Warumungu Body Book}, a translation booklet by the Media Unit. A 2004 ANU newsletter mentions it alongside Warumungu language programmes at the high school and, more recently, the primary school: \url{https://www0.anu.edu.au/linguistics/nash/aust/pak/2.1-PAKnewsletter-Apr2004.pdf}.
\end{itemize}

\section{Linguistic and Reference Materials}
\begin{itemize}
\item James Simpson, \emph{A Learner's Guide to Warumungu} (IAD Press, 2002): \url{https://www.iadpress.com.au/products/a-learners-guide-to-warumungu}.
\item The AIATSIS AUSTLANG entry C18 lists historical sources, including Wafer's 1982 picture dictionary and the 2002 learner's guide: \url{https://collection.aiatsis.gov.au/austlang/language/c18}.
\item \emph{Oral History of Jimmy Jampin Jones} (2000), a booklet compiled in a two-week workshop at Batchelor and meant for a community language centre. It is listed under the title ``Warumungu Language Resources''. I saw no direct URL for it.
\end{itemize}

\section{Religious and Community Resources}
Papulu Apparr-kari Aboriginal Corporation in Tennant Creek is the language centre named in results. A 2004 newsletter from it is cited above. I saw no scripture resources and give no contact details.

\section{Digital and Audio Resources}
I found no standalone audio or digital resources specific to Warumungu.

\section{Curriculum Development Context}
Warumungu is not recorded as a regular school subject. The 2004 newsletter, however, mentions language programmes at the Tennant Creek high school and primary school. The dictionaries, the picture dictionary and the body book offer a starting vocabulary base. The learner's guide gives a grammar and pronunciation foundation. The available print materials are dated, so curriculum work should be done with the regional language centre and speakers.
\end{document}